%%%%%%%%%%%%%%%%%%%%%%%%%%%%%%%%%%%%%%%%%%%%%%%%%%%%%%%%%%%%%%%%%%%%
% - Desempenho dos modelos supervisionados
% - Comparação entre modelos discriminativos e espaço latente
% - Impacto da integração dos dados clínicos
% - Resultado da abordagem alternativa
% - Interpretação biológica
% - Discussão limitações
%%%%%%%%%%%%%%%%%%%%%%%%%%%%%%%%%%%%%%%%%%%%%%%%%%%%%%%%%%%%%%%%%%%%

\chapter{Conclusions and Discussion}
\label{ch:discussion}
 
% ================================================================
\section{Overview}
\label{sec:disc_overview}
% ================================================================
 
This chapter synthesises the findings of this dissertation and addresses three
central questions that motivated the work: what are the main conclusions; whether
the research hypotheses formulated in Chapter~\ref{cha:introduction} have been
answered; and in what ways the results obtained can bring value to the real world.
 
The experimental work was conducted in two phases. A \textbf{pilot study}
(Chapter~\ref{ch:pilot}) developed, validated, and tuned the analytical pipeline on a
244-patient cohort curated by our partners at Nova Medical School, establishing a
methodological foundation and a first set of biological findings. The
\textbf{main analysis} (Chapter~\ref{ch:expanded}) re-applied the same pipeline, 
without any modification, to an independently assembled, larger cohort of 702 patients
drawn directly from cBioPortal \cite{cerami2012cbio, TCGA2012}, testing whether the pilot findings generalise at
scale. The results of both phases are brought together here to draw final conclusions.
 
Three complementary analytical approaches were used throughout both phases: a
supervised discriminative framework (XGBoost under a One-vs-Rest formulation),
a statistical validation layer (ANOVA F-test feature ranking), and an unsupervised
geometric analysis (K-Means clustering). Because the three methods make no shared
modelling assumptions, any conclusion that emerges consistently from all three carries
a level of evidential weight that no single method could achieve alone. The degree of
convergence (or divergence) between them is therefore central to the discussion that
follows.
 
% ================================================================
\section{The Separability Gradient: A Consistent Biological Finding}
\label{sec:disc_gradient}
% ================================================================
 
The most robust and replicable result of this dissertation is the existence of a
\textbf{consistent separability gradient} across the four PAM50 \gls{bc}
subtypes in \gls{mirna} expression space:
 
\begin{equation}
    \text{Basal-like} \;\succ\; \text{Luminal A} \;\succ\;
    \text{HER2-Enriched} \;\approx\; \text{Luminal B}
    \label{eq:gradient}
\end{equation}
 
This ordering was first observed in the pilot cohort, where Basal-like was
classified with perfect F1-score (1.00), Luminal A with moderate performance
(F1 = 0.81), and both HER2-Enriched and Luminal B at substantially lower performance
(F1 = 0.50 each) under miRNA-only conditions, where this was subsequently confirmed in the
expanded cohort with near-identical ordering (Basal: F1 = 0.98, Luminal A:
F1 = 0.89, HER2-Enriched: F1 = 0.42, Luminal B: F1 = 0.56) - these results converge with 
the one seen in the literature \cite{GAIN_BRCA_Patel2025}. Crucially, this
gradient was independently corroborated by both the ANOVA F-test sweep and the
K-Means clustering purity analysis in both cohorts, establishing that it is not a
model-specific or cohort-specific artefact, but a \textbf{genuine geometric property
of the miRNA expression space}.
 
The consistency of this gradient across two independently assembled datasets, 
differing in size (244 vs.\ 702 samples), clinical annotations (receptor status
vs.\ hypoxia and genomic instability scores), miRNA nomenclature conventions, and
preprocessing histories, constitutes a strong form of experimental replication. It
strongly suggests that the gradient reflects \textbf{a real biological property of
PAM50 subtypes as encoded in miRNA expression}: Basal-like tumours possess a distinct
and compact miRNA programme; Luminal A has a detectable but more heterogeneous
signature; and HER2-Enriched and Luminal B are transcriptomically proximal, their
miRNA profiles lacking the geometric separability required for reliable algorithmic
recovery. This conclusion is consistent with the well-established clinical
distinctiveness of Basal-like (triple-negative) breast cancer
\cite{bc_molecular_Perou2000} and with the known transcriptomic proximity between
the two Luminal subtypes.
 
Importantly, this finding was first established in the pilot cohort (Chapter~\ref{ch:pilot})
and then confirmed in the expanded cohort (Chapter~\ref{ch:expanded}). The fact that
the pilot correctly anticipated
the ordering and approximate magnitudes of separability seen in the expanded analysis
validates the pilot study as a reliable methodological scaffold, not merely a
preliminary exercise.
 
% ================================================================
\section{Cross-Paradigm Validation and the Proposed Feature Panels}
\label{sec:disc_panels}
% ================================================================
 
Beyond the separability gradient, the central biological question of this dissertation
concerns the \textbf{identity of the features that drive the separation}. Three
independent analytical approaches - the XGBoost booster gain importances, the ANOVA
F-test ranking, and the K-Means centroid Z-scores - each produce a ranked list of
candidate features per subtype. Features that appear consistently across all three
methods represent the strongest biomarker candidates, as their discriminative power
cannot be attributed to the assumptions or biases of any single method. This section
presents the proposed panels derived from the expanded cohort and then assesses their
cross-cohort robustness.
 
% ---------------------------------------------------------------
\subsection{The Expanded Cohort Biomarker Panel}
\label{subsec:disc_expanded_panel}
% ---------------------------------------------------------------

The proposed panel is not uniformly strong across subtypes; the following paragraphs
examine each subtype in turn.
 
\paragraph{Basal-like: the most robust and actionable panel.}
The Basal-like panel is the most compact and most multiply validated of the four.
Six miRNAs : \textbf{MIR-577, MIR-934, MIR-18a, MIR-135b, MIR-671, and MIR-190b},
 are doubly validated by both XGBoost and ANOVA in the expanded cohort, with
MIR-190b additionally appearing in the Luminal B panel (see the intersection in
Figure~\ref{fig:venn_expanded}). Of these, \textbf{MIR-577} is by far the most
dominant feature: it accounts for 46.6\% of the total XGBoost booster gain under
miRNA-only conditions and for 64.2\% when clinical features are added. MIR-577 and
MIR-18a are additionally corroborated by the K-Means centroid Z-score analysis,
making them \textbf{triply validated} candidates. The concentration of discriminative
power in a small set of miRNAs, combined with a near-perfect supervised classification
performance (F1 = 0.98) and a 95.5\% unsupervised cluster purity, makes Basal-like
the most biologically credible and practically actionable case for miRNA-based subtype
identification in this work.
 
\paragraph{Luminal A: a learnable boundary, an internally heterogeneous subtype.}
Luminal A presents a more nuanced picture. Two miRNAs: \textbf{MIR-590} and
\textbf{MIR-130b}, are doubly validated in the miRNA-only condition, with MIR-590
being the leading feature by XGBoost gain (8.84\% of total gain). The Ragnum Hypoxia
Score is the top-ranked clinical feature when clinical annotations are added, reflecting
the annotation-specific nature of the clinical signal available in the expanded cohort.
However, the unsupervised analysis consistently fragments Luminal A into two distinct
sub-clusters with non-overlapping centroid signatures, suggesting that this subtype
harbours genuine \textbf{internal heterogeneity} that the PAM50 classification does
not resolve. The practical implication is that a single Luminal A miRNA panel may be
insufficient; a deployed diagnostic panel would need to account for this intra-subtype
structure.
 
\paragraph{HER2-Enriched: a diffuse panel requiring clinical complementation.}
For HER2-Enriched, the doubly-validated miRNA panel consists of \textbf{MIR-4728,
MIR-2115, MIR-30a,} and \textbf{MIR-6733}. The Buffa Hypoxia Score is additionally
validated as a clinical discriminator in the miRNA + clinical condition, consistent
with the known association between HER2 amplification and tumour hypoxia
\cite{bc_subtypes_Prat2015Clinical}. Classification performance, however, remains
limited (F1 = 0.42 miRNA-only; F1 = 0.57 with clinical features) and no high-purity
unsupervised cluster was recovered for this subtype in either cohort. Both paradigms,
across both datasets, converge on the same conclusion: \textbf{miRNA expression alone
is insufficient to define a robust discriminative panel for HER2-Enriched tumours.}
Any practical diagnostic panel for this subtype would require the integration of
additional molecular data, receptor status or hypoxia profiling, alongside miRNA
expression.
 
\paragraph{Luminal B: the most analytically challenging subtype.}
Luminal B is the most challenging subtype across all methods. The doubly-validated
miRNA panel in the expanded cohort is limited to \textbf{MIR-193b} and
\textbf{MIR-301a} (miRNA-only), with MIR-190b added in the miRNA + clinical
condition; this last feature, however, is shared with the Basal-like panel
(Figure~\ref{fig:venn_expanded}). Luminal B's transcriptomic proximity to Luminal A,
 confirmed geometrically by the K-Means analysis in both cohorts, limits the
identifiability of a compact, subtype-specific miRNA signature. The improvement in
F1-score from 0.56 to 0.65 with clinical features (in particular Ragnum Hypoxia
Score and MSI-related metrics) is meaningful but not sufficient to fully resolve
the Luminal A / Luminal B boundary.

Figure~\ref{fig:venn_expanded} summarises the proposed biomarker panel from the
expanded cohort analysis. The features shown in each subtype region are
\textit{doubly validated} by both XGBoost gain importances and ANOVA F-test ranking.
Features additionally corroborated by the K-Means centroid Z-scores - making them
\textit{triply validated} - are marked with~$\ast$.
 

\begin{figure}[!ht]
    \centering
    \begin{subfigure}{0.95\textwidth}
        \centering
        \includegraphics[width=0.90\textwidth]{diagrama_venn_crosscohort_mirna.drawio.png}
        \caption{miRNA only}
    \end{subfigure}
    \vspace{0.5cm}
    \begin{subfigure}{0.95\textwidth}
        \centering
        \includegraphics[width=0.90\textwidth]{diagrama_venn_crosscohort_mirna+clinical.drawio.drawio.png}
        \caption{miRNA and Clinical data}
    \end{subfigure}

    {\footnotesize \textbf{Note:} The prefix "MIR-" is omitted for brevity and to improve readability.}
    
    \caption{Feature intersections showing doubly-validated markers for all subtypes under both miRNA-only and miRNA+clinical conditions.}
    \label{fig:venn_expanded}
\end{figure} 
 
% ---------------------------------------------------------------
\subsection{Cross-Cohort Validated Panel}
\label{subsec:disc_cross_panel}
% ---------------------------------------------------------------
 
The pilot and expanded cohorts differ in their preprocessing pipelines, clinical
annotations, miRNA nomenclatures, and cohort sizes. Cross-cohort replication -
features that appear as doubly-validated biomarker candidates \textit{in both}
cohorts - therefore represents the strongest validation standard in this work: a
finding that is stable across all sources of systematic variation between the two
datasets.
 
Figure~\ref{fig:venn_crosscohort} presents the cross-cohort validated panel. Only
features identified as doubly-validated candidates in both the pilot cohort
($n = 244$) and the expanded cohort ($n = 702$) are shown. Features that appeared
only in the pilot cohort are annotated in red, the features that appeared only in 
the expanded cohort are annotated in orange, while the features that appeared in 
both cohorts are annotated in \textbf{bold}.
 
\begin{figure}[h!]
    \centering
    % \includegraphics[width=0.90\textwidth]{venn_crosscohort_panel.pdf}
    \caption{Cross-cohort validated miRNA panel. Only features identified as
    doubly-validated biomarker candidates in \textit{both} the pilot cohort
    ($n = 244$) and the expanded cohort ($n = 702$) are shown. Dashed ellipses
    represent subtypes for which no cross-cohort replicated panel was established.
    The Basal-like subtype is the only one with a cross-cohort validated panel,
    comprising four miRNAs (MIR-577, MIR-934, MIR-18a, MIR-135b). Dashed callouts
    indicate miRNAs exclusive to the pilot (MIR-505) or the expanded (MIR-190b,
    MIR-671) cohort.}
    \label{fig:venn_crosscohort}
\end{figure}
 
The result is unambiguous: \textbf{only the Basal-like subtype yields a cross-cohort
validated miRNA panel}. Four miRNAs: \textbf{MIR-577, MIR-934, MIR-190B} and
\textbf{MIR-135b} appear as top-ranked features in both the XGBoost gain
importances and the ANOVA F-test ranking for Basal-like, independently, in both
datasets. These four miRNAs are therefore supported by three converging layers of
evidence: multivariate discriminative power (XGBoost), marginal statistical
separation (ANOVA), and cross-cohort replication.
 
For the remaining subtypes, the absence of cross-cohort validated features is itself
an informative finding. It is not a power problem: the expanded cohort contains almost
three times as many samples as the pilot and the separability gradient is preserved.
Rather, it reflects the biological and statistical nature of these subtypes: their
miRNA signatures are more diffuse, more context-dependent (as evidenced by the
annotation-dependence of the clinical features that enter their panels), and
insufficiently concentrated to produce a stable cross-dataset biomarker set. This
conclusion is consistent with the broader literature on Luminal B and HER2-Enriched
breast cancer, where robust and reproducible miRNA panels have remained elusive.
 
% ================================================================
\section{Revisiting the Research Hypotheses}
\label{sec:disc_hypothesis}
% ================================================================
 
Chapter~\ref{cha:introduction} articulated several interconnected research hypotheses.
This section revisits each one in light of the results obtained.
 
\paragraph{Can miRNA expression profiles discriminate between breast cancer subtypes?}
\textbf{Partially, and in a structured way.}
MiRNA expression profiles can reliably discriminate Basal-like tumours from all others,
and Luminal A to a meaningful but lesser extent. For HER2-Enriched and Luminal B, the
miRNA expression space does not provide a compact enough geometric separation for
reliable discrimination by either supervised or unsupervised means. This answer is
itself scientifically valuable: it rules out both the pessimistic hypothesis (no
signal) and the optimistic hypothesis (complete signal), clarifying the true scope and
limits of miRNA-based subtyping.
 
\paragraph{Can machine learning identify a compact panel of discriminative miRNA biomarkers?}
\textbf{Yes, for Basal-like; tentatively for others.}
The XGBoost and ANOVA approaches consistently identify a small set of miRNAs dominated
by MIR-577 for Basal-like, accounting for the vast majority of the discriminative
signal. For the other subtypes, the signal is distributed across many features and a
compact panel of sufficient clinical relevance has not been conclusively established.
 
\paragraph{Does unsupervised learning reveal implicit structure in the miRNA expression space?}
\textbf{Yes, and it adds independent scientific value.}
K-Means, operating without any label supervision, recovers a Basal-like cluster with
up to 100\% purity in the expanded cohort, corroborating the same miRNAs identified
by the supervised approach. Beyond validation, the unsupervised analysis provides
observations that the supervised approach alone could not: it reveals the internal
heterogeneity of Luminal A (consistently split into two sub-clusters) and places the
difficulty of HER2-Enriched and Luminal B in a geometric context, confirming that
their inseparability is a data-intrinsic property rather than a model limitation.
 
\paragraph{Do clinical features improve classification?}
\textbf{Yes, selectively.}
Clinical features improve the supervised classification of HER2-Enriched
($+0.15$ F1) and Luminal B ($+0.09$ F1) but leave Basal-like essentially unchanged,
confirming that the miRNA signal for Basal-like is self-sufficient. Importantly,
however, the clinical features most useful in the pilot (receptor status: ER, PR,
HER2) are not the same as those in the expanded cohort (Buffa and Ragnum Hypoxia
Scores, MSI metrics), highlighting that the identity of useful clinical covariates is
\textit{annotation-dependent} and may not generalise across institutional data
releases.
 
\paragraph{Do the findings generalise across independent datasets?}
\textbf{Yes for the separability gradient and the Basal-like panel; not for the
specific feature panels of the other subtypes.}
The separability ordering (Basal $>$ Luminal A $>$ HER2 $\approx$ Luminal B) is
reproduced with near-identical rankings in both cohorts. The Basal-like feature
panel (MIR-577, MIR-934, MIR-18a, MIR-135b) is cross-cohort replicated. The panels
for Luminal A, HER2-Enriched, and Luminal B should be treated as
\textbf{hypothesis-generating} observations specific to their respective cohorts'
annotations and sample compositions.
 
% ================================================================
\section{Real-World Applicability and Future Directions}
\label{sec:disc_applications}
% ================================================================
 
The computational work developed in this dissertation does not exist in isolation
from experimental biology. As stated in Chapter~\ref{cha:introduction}, the findings
of this work could be tested through \textit{in vitro} and \textit{in vivo} experiments
conducted by the research team of Dr.\ Bárbara Mendes at NOVA Medical School
\cite{_Mendes2022Nanodelivery}. This direct collaboration between computational
biomarker discovery and laboratory biology represents a critical bridge between data
analysis and clinical relevance.
 
The cross-cohort replicated panel, \textbf{MIR-577, MIR-934, MIR-18a, and
MIR-135b}, is the most well-positioned target for this biological validation.
The fact that these miRNAs are doubly validated computationally \textit{and}
cross-cohort replicated substantially reduces the risk of false positives compared
to findings derived from a single dataset or a single method. If laboratory
validation confirms their differential expression and functional relevance to
Basal-like tumour biology, these miRNAs could form the basis of a targeted
diagnostic panel for the identification of triple-negative breast cancer - a
subtype for which targeted molecular biomarkers are particularly needed, given the
absence of hormone receptor targets that would otherwise guide therapy.

 
% ---------------------------------------------------------------
\subsection{Limitations}
\label{subsec:disc_limitations}
% ---------------------------------------------------------------
 
Several limitations must be acknowledged to contextualise the findings appropriately.
 
\textbf{Single-study origin.} Both datasets originate from TCGA-BRCA, a single large
multi-omics study conducted under a uniform experimental protocol. Although the two
releases are independently processed and their patient sets partially non-overlapping,
they cannot be considered fully independent biological replicates. External validation
on entirely distinct cohorts, such as METABRIC \cite{bc_molecular_Perou2000} or
publicly available GEO datasets, would be required to establish the generalisability
of the proposed panels beyond the TCGA ecosystem.
 
\textbf{Absence of a healthy control set.} As discussed in
Section~\ref{sec:challenges+research-hypothesis}, neither dataset includes miRNA
expression profiles from healthy breast tissue. The models therefore discriminate
subtypes from one another, but cannot determine whether a given miRNA signature is
cancer-specific. A control group would not only expand the scope of the classification
task but also improve the specificity of the identified biomarker, as demonstrated
in related work on gastric cancer \cite{ml_gastric_Azari2023}.
 
% ---------------------------------------------------------------
\subsection{Future Directions}
\label{subsec:disc_future}
% ---------------------------------------------------------------
 
The findings of this dissertation open several promising directions for future research.
 
The most immediate next step is the biological validation by the NOVA Medical School
team. If the cross-cohort replicated Basal-like panel is confirmed experimentally,
the logical subsequent step is the development of a \textbf{minimal clinical panel}
by identifying the fewest miRNAs from the proposed set that achieve sufficient
classification accuracy in a targeted assay context.
 
The internal heterogeneity of Luminal A identified by the K-Means analysis warrants
dedicated investigation. A sub-clustering analysis with varying values of $K$,
combined with clinical outcome data (survival, treatment response), could determine
whether the two Luminal A sub-populations correspond to clinically distinct entities
with different prognoses or treatment sensitivities.
 
For HER2-Enriched and Luminal B, the integration of additional omics layers alongside
miRNA data (such as mRNA expression, copy number variation, or DNA methylation)
might provide the orthogonal signal necessary to build compact, reliable discriminative
panels. These multi omnic data sources could reveal new relations in the biological 
dynamics of the human body and the nature of breast cancer subtyping.
 
% ================================================================
\section{Summary}
\label{sec:disc_summary}
% ================================================================
 
This dissertation investigated the applicability of machine learning - supervised and
unsupervised - to the classification of breast cancer subtypes using microRNA
expression and clinical data from two independent releases of the TCGA-BRCA cohort.
The main conclusions are as follows.
 
\textbf{First}, a consistent separability gradient was identified and replicated across
both cohorts: \mbox{Basal-like $>$ Luminal A $>$ HER2-Enriched $\approx$ Luminal B}.
This gradient is a genuine, dataset-independent property of the PAM50 miRNA expression
space.
 
\textbf{Second}, a cross-paradigm and cross-cohort validated biomarker panel was
established for the Basal-like subtype, comprising \textbf{MIR-577, MIR-934, MIR-18a,
and MIR-135b}, with MIR-577 exhibiting the strongest and most consistent discriminative
signal across all methods and both datasets. For the other subtypes, the proposed
panels are hypothesis-generating rather than conclusive.
 
\textbf{Third}, the unsupervised K-Means analysis provided independent corroboration
of the supervised findings for Basal-like and added novel insight into the internal
heterogeneity of Luminal A: a dimension of the data structure that the supervised
approach alone could not have revealed.
 
\textbf{Fourth}, clinical features carry discriminative signal that supplements miRNA
expression for HER2-Enriched and Luminal B, but the identity of the most useful
clinical covariates is annotation-dependent, limiting cross-cohort generalisability
of any clinical-feature-dependent panel.
 
These findings are ready to be tested biologically through the ongoing collaboration
with NOVA Medical School, and they lay a computational foundation for the development
of targeted, evidence-based miRNA panels for breast cancer subtype stratification,
with the Basal-like panel representing the most mature and actionable candidate.

\section{Future Work}
\label{sec:future_work}
% ================================================================

While this work established a robust computational baseline, several avenues
remain to deepen the biological insights and predictive accuracy of breast
cancer subtyping.

% ---------------------------------------------------------------
\paragraph{Multi-Omic Integration}
\label{subsec:future_multiomics}
% ---------------------------------------------------------------

A primary limitation of the present study is the difficulty in separating
Luminal B and HER2-Enriched subtypes using miRNA expression alone. Future
research should transition towards a \textbf{multi-omic approach} that
combines miRNA with complementary data modalities. Integrating epigenetic
layers such as DNA methylation or copy number variation (CNV) could provide
the orthogonal signal necessary to resolve the boundaries where miRNA
signatures overlap. Beyond genomics, the incorporation of protein expression
data would allow for the direct validation of miRNA-target interactions,
bridging the gap between transcriptomic signatures and the functional
biological mechanics that drive subtype-specific tumour behaviour.

% ---------------------------------------------------------------
\paragraph{Advanced Algorithmic Architectures}
% ---------------------------------------------------------------

As biological datasets continue to scale and encompass more diverse data
modalities, the limitations of classical ML may be effectively addressed by
more complex models. Deep learning architectures (such as Autoencoders for
non-linear dimensionality reduction or Graph Neural Networks for modelling
molecular interaction networks) could capture dependencies that linear or
tree-based classifiers miss. Simultaneously, evolutionary and genetic algorithms
for feature selection could further refine the search for minimal clinical
panels, identifying the most potent biomarker combinations across
high-dimensional multi-omic spaces without relying on predefined assumptions
about feature relevance.

% ---------------------------------------------------------------
\paragraph{Control Group}
% ---------------------------------------------------------------

As a limitation of the present study, neither dataset includes miRNA expression 
profiles from healthy breast tissue. The models therefore discriminate subtypes 
from one another, but cannot determine whether a given miRNA signature is 
cancer-specific. A control group would not only expand the scope of the 
classification task but also improve the specificity of the identified 
biomarker, as demonstrated in related work on gastric cancer 
\cite{ml_gastric_Azari2023}.


% ---------------------------------------------------------------
\paragraph{Experimental Validation}
% ---------------------------------------------------------------

The computational findings of this dissertation do not exist in isolation
from experimental biology. The cross-cohort validated Basal-like panel
(MIR-577, MIR-934, MIR-18a, and MIR-135b) is currently positioned for
\textit{in vitro} and \textit{in vivo} testing at NOVA Medical School, in
collaboration with the research team of Dr.\ Bárbara Mendes~\cite{_Mendes2022Nanodelivery}.
Future studies should prioritise these functional assays to confirm the
biological relevance of the identified miRNAs and to establish the minimum
effective panel size for a reliable diagnostic readout. If validated
experimentally, these candidates could form the basis of a targeted
diagnostic tool for Basal-like (triple-negative) breast cancer, a subtype
for which robust molecular biomarkers are particularly needed, given the
absence of hormone receptor targets that would otherwise guide therapy.
